\section{The Record as a Plausibility Probe}\label{sec:evidence}

The parameters of the model are not measured, and a negative exit premium can rationalize almost any act of compliance after the fact. Compliance alone also cannot separate coercion from shared interest, because successful coercion ends at the threat stage \citep{drezner2003hidden}. We therefore ask what the model implies that its rivals do not, fix rules for coding what Washington gave and threatened, and read the Dutch record and nine comparison cases with those rules. None of this is a test. \Cref{sec:tests} describes the designs that would be.

\subsection{Observable implications}\label{sec:predictions}

\Cref{tab:implications} sets seven implications of the model beside five rival accounts of allied compliance. Under shared threat perception the Netherlands aligns because it shares Washington's view of China, which in the model is a large intrinsic cost of exit $\eta$. In \citet{lake2009hierarchy} compliance is owed to a patron whose authority rests on protection already provided and partly on legitimacy; in \citet{morrow1991alliances} the protégé trades autonomy for security, so protection is itself the payment. \citet{drezner1999sanctions} argues that adversaries, who expect future conflict, rarely win concessions by coercion, while allies win them when they try. A cost-minimizing coercer polices each demand with the cheapest credible instrument and sizes sanctions so that punishing a small offence does not erode deterrence of a large one \citep{stigler1970optimum, shavell1992note}.

\begin{table}[t]
\centering
\scriptsize
\setlength{\tabcolsep}{3pt}
\caption{Observable implications of the model and of rival accounts}\label{tab:implications}
\begin{tabularx}{\textwidth}{@{}>{\raggedright\arraybackslash}p{2.35cm}*{6}{>{\raggedright\arraybackslash}X}@{}}
\toprule
Implication (result) & This model & Shared threat (large $\eta$) & Hierarchy & Autonomy--security trade-off & Conflict expectations & Cost minimization\\
\midrule
Threshold (\cref{cor:master}) & Holder is paid only if $Q>0$, or if talks risk rupture and it has agenda power & No payment; holder aligns for its own reasons & No payment; compliance is owed & Protection is the payment & Allies concede without payment & No prediction about payment\\
\addlinespace
Linkage only with slack (\cref{prop:counter}(b)) & Linkage against secondary demands appears when the holder has slack and fades as $Q\to0$ & Little coercion of any kind & Commands need no threats & Protection is linked to autonomy throughout & Coercion of allies rare but effective, whatever the rival offers & Cheapest credible instrument, unchanged as $Q$ moves\\
\addlinespace
One budget across issue areas (\cref{prop:slack}(c), \cref{cor:rationing}) & Demands on trade, defence and controls are substitutes; the budget shrinks as the rival's offer rises & Issue areas independent & Authority bounded issue by issue by legitimacy & One autonomy price for security, set within the alliance & No cross-issue constraint & Instruments priced separately\\
\addlinespace
Counter-threats (\cref{prop:counter}(a)) & Change the patron's demand mix and tax the patron; executed only to crowd out other demands & Holder absorbs them and keeps aligning & Patron shields its subordinate from third parties & No prediction & Rival's coercion of an adversary yields little & No prediction\\
\addlinespace
Rival bids (\cref{prop:conditionality}(c)(i)) & A stronger bid elicits a more conditional guarantee, not a payment, while assurance is cheap & No response needed & Patron must provide more order to keep its authority & Better terms for the protégé & No prediction & No prediction\\
\addlinespace
Preventive control (\cref{prop:control}, \cref{cor:windows}) & Control is demanded while the holder is weak but expected to strengthen & No demand & More control where dependence is deeper, with no window & No prediction & No prediction & Control demanded whenever cheapest, with no window\\
\addlinespace
Frontline states (\cref{rem:council}) & Members with large conditional exposure join blocking coalitions at no cost, though they cannot block alone & Members most hostile to China and Russia block & Most dependent members block & Members receiving most security block & No prediction & No prediction\\
\bottomrule
\end{tabularx}
\par\smallskip
\begin{minipage}{\textwidth}\scriptsize
\emph{Note.} Entries for the rival accounts are our derivations from \citet{lake2009hierarchy}, \citet{morrow1991alliances}, \citet{drezner1999sanctions} and \citet{stigler1970optimum}, not claims their authors make about these cases. $Q$ is the exit premium (\cref{def:Q}).
\end{minipage}
\end{table}

The threshold separates the model from hierarchy and the autonomy--security trade-off only through holders with a positive premium, which should be paid, as Kyrgyzstan was for Manas when a Russian aid package raised the rent Washington paid \citep{cooley2012great}. The second and third rows discriminate most. In both, what Washington can obtain shrinks as China's credible offer rises, whatever its instruments cost. Cost minimization implies no such change, and none of the alliance accounts lets a rival's offer enter the price of compliance. In the fourth and fifth rows, hierarchy implies a patron that shields its subordinate, while our model implies that a Chinese counter-threat changes what Washington asks for and that a stronger Chinese bid makes the guarantee more conditional. The frontline row discriminates least: every account in which security dependence produces deference shares it.

\subsection{Coding transfers, threats and leverage}\label{sec:coding}

Reading the record through \cref{prop:master} requires a line between a payment and a threat that was never carried out. We draw it at Washington's ex post optimal provision to an aligned ally, the baseline $x_j(A)$ of each instrument (\cref{sec:security}).

A \emph{transfer} is a benefit conditional on alignment that exceeds what Washington would provide to an aligned ally if it optimized ex post. Only transfers enter $P^*$. A \emph{threat left unexercised} is either a baseline benefit that Washington would withdraw after exit or a sanction it could impose and does not. It enters the exit cost $c$ as a conditional part $d_j$, however valuable it is to the holder. \Cref{rem:gatekeeper} draws the same line: a benefit Washington provides to a loyal Netherlands anyway, and withdraws on exit, lowers the gatekeeper's premium, while a promise above the baseline is a payment like any other.

The December 2024 exclusion of the Netherlands and the other countries on Supplement No.~4 from the new SME licence requirements (\cref{sec:background}) is the case that matters. BIS conditioned it on the listed governments' own capacity: they ``have the authority to control key SME items of concern'' \citep{bis2024fdpr}. The exclusion lasts as long as The Hague controls the items itself, and BIS can withdraw it by amending the rule. We therefore code it as a threat left unexercised, conditioned on Europe's own capacity to control, or equivalently as a reward for self-binding (\cref{rem:bind}). It is not a transfer. With this coding the exclusion is consistent with a nonpositive premium: it adds to $c$ and lowers $Q$, and under \cref{cor:master} a holder with $Q\le 0$ and little risk of rupture receives no transfer. The same rule codes the ``more favorable treatment'' that the National Security Strategy offers countries aligning their export controls as conditionality (\cref{sec:conditionality}).

We code the instrument a US demand invokes in three classes. \emph{Flow jurisdiction} reaches China-bound flows directly (de minimis and foreign direct product rules, entity listings, licences for US-origin items); it leaves $Q$ unchanged and needs no concession from the holder (\cref{prop:counter}(b)). \emph{Conditional linkage} threatens to withdraw from an aligned holder something in another domain that an exited holder would lose anyway, such as posture, AI access or trade terms. An \emph{unconditional} instrument hurts the holder whatever it does with the chokepoint. Chinese responses are coded as counter-threats if announced before the holder concedes and as executed if carried out.

Transfers are a narrow measure of leverage, on which the Netherlands has little to show almost by construction. Four broader measures also have meaning in the model: \emph{delay} between a US move and the holder's step, the \emph{scope} of the holder's controls relative to the demand, \emph{retained jurisdiction} over licensing, and \emph{tolerated hedges}, flows Washington could reach but leaves alone. None is a payment. On our reading they measure how much of the holder's slack Washington chose to spend: by \cref{cor:rationing} a patron leaves a hedge alone when its value per unit of the holder's cost is low, and retained jurisdiction is self-binding with the key kept in Europe (\cref{rem:bind}).

\subsection{The Dutch record and Nexperia}\label{sec:record}

\paragraph{Alignment without a price.} Since 2019 The Hague has aligned in steps, and from 2023 each step followed a US move (\cref{sec:background}). Under our rules it received no transfer: its one benefit, Supplement No.~4, is a threat left unexercised. If talks carry no risk of rupture, or Washington makes the offers, \cref{cor:master}(b) gives $P^*\le 0$ whenever $Q\le0$, and the record is consistent with a holder whose premium has been negative throughout, even though customers in China took between 26 and 36 percent of ASML's sales from 2023 to 2025 (\cref{sec:background}).

The negative premium needs no security term. On our reading, US flow jurisdiction over China-bound tools blocks most of what China would gain from exit, which makes $\theta$ large and $\omega$ small (\cref{eq:omega}). ASML's US entities and staff are US persons (\cref{sec:background}), so the firm would not ship against a credible US penalty even if The Hague licensed exit, and Europe would then have to compel it (\cref{rem:firm}). With $\kappa=0$ these terms alone can make $Q$ negative. The record therefore supports the jurisdictional channel and leaves the security channel untested: we have found no documented case that links US security commitments to ASML. Cross-domain linkage appears elsewhere (\cref{sec:cases}), never on Dutch chip controls. The disputes since 2023 concern legacy sales and servicing, a hedge that by \cref{prop:counter}(b) Washington polices with flow jurisdiction whenever that is credible, so they could not reveal $\kappa$ even if it were large.

Nor do the security signals of 2025 and 2026 settle it. An untargeted drawdown should raise Europe's premium, and a withdrawal aimed at an ally's conduct is an executed threat (\cref{sec:conditionality}), yet the Dutch response to US pressure in 2026 was voice without exit. Either jurisdiction alone kept the premium negative, or perceived conditionality rose with the targeted moves. The record cannot tell these apart.

\paragraph{Slack spent and unspent.} The broader measures show what The Hague kept. Licensing of the older immersion tools came a year after the advanced ones, The Hague took the licensing pen from Washington in September 2024, and servicing in China continued at least into mid-2026 (\cref{sec:background}). We read this as rationing: Washington spent Europe's slack first on the sales controls it valued most and left the servicing hedge, costly to ASML, until 2026. A shared-interest account fits the substance of Dutch alignment but not its timing, since a holder aligning for its own reasons need not wait for US rules.

\paragraph{Nexperia.} The Nexperia dispute of 2025 is the one episode in which a Chinese counter-threat was executed against a Dutch step taken under US pressure. On 5 June 2025 the Dutch economy ministry told Nexperia that Washington might soon introduce a ``50\% rule'' and that, if the Nijmegen firm sought an exemption, Washington would look at measures to stop transfers of US intellectual property and technology. On 12 June the Dutch Foreign Ministry met the State Department's Bureau of International Security and Nonproliferation. The Dutch minutes record that the US side called it problematic that the chief executive was still the same Chinese owner, and conclude that it was ``almost certain the CEO will have to be replaced to qualify for the exemption from the entity list'' \citep[paras.~2.19--2.20; our translation]{gerechtshof2025nexperia}. The minutes, reported in October 2025 \citep{dutchnews2025nexperia}, are a Dutch record of US views, not a formal US condition. On 30 September, after a US rule threatened Nexperia's operations, the economy minister used the Goods Availability Act to take control of the firm, and the Enterprise Chamber suspended the chief executive in early October \citep{taylorwessing2025nexperia, gerechtshof2025nexperia}. China then blocked exports from Nexperia's Chinese plants, which threatened production at European carmakers \citep{taylorwessing2025nexperia}. At Busan on 30 October, Washington agreed to suspend the Affiliates Rule for one year from 10 November, and China agreed to ``take appropriate measures to ensure the resumption of trade from Nexperia's facilities in China'' \citep{fmprc2025busan, whitehouse2025busan}. On 19 November the Dutch minister suspended his order \citep{taylorwessing2025nexperia}. The Enterprise Chamber's measures were separate from the order and outlasted it. % TODO(bib): source is Kamerstukken 32637 nr. 717 (key kst32637_717 in notes/review/gap_dutch_domestic_politics.md), not in refs.bib.

We read the episode as a US-backed demand on the Netherlands (an entity-list exemption in exchange for removing Chinese control of the firm), met by a Chinese counter-threat that was executed, with the Dutch suspension following Washington's suspension of the rule behind the demand. Under complete information, \cref{prop:counter}(a)(iii) allows execution on the equilibrium path only when the counter-threat crowds out another demand. That reading requires three things: a Chinese threat known in The Hague and Washington before 30 September; a US demand on the Netherlands that Washington shelved because the Nexperia concession had become expensive; and a Dutch government that conceded while expecting to bear the cost. The alternative, outside our complete-information model, is that The Hague, and perhaps Washington, misjudged China's resolve, which requires that the Chinese response surprised them. The reversal does not discriminate, since by then the rule behind the US demand had been suspended. Neither reading is established by the public record, and the minutes that revealed the US demand exist only because a court published them.

Europe bore the cost and was not compensated. By \cref{prop:counter}(a)(v) a counter-threat leaves the holder's loss unchanged when Washington has a divisible sunk demand that China cannot reach, and otherwise the holder's loss can rise. The relief came through a US concession to China at Busan. US--China bargaining is outside the model, so we do not read Busan as a transfer to Europe. The reversal under Chinese pressure is also evidence that compelling a firm is costly when China retaliates against the compulsion itself (\cref{rem:firm}). And the order shows that the Dutch state can override control of a firm under its own law, which on our reading lowers the enforceability of any control right that Washington might hold through a Dutch entity (\cref{cor:control}(c)).

\paragraph{Servicing and the MATCH Act.} Servicing of China's installed base is the hedge in our terms, and Washington tolerated it into 2026. The MATCH Act would police it by flow jurisdiction: allies that do not align within 150 days would face US controls on sales and servicing (\cref{sec:background}). In the model a statute is a stake that makes the threat credible if $e_F\le\bar m_U$, and a credible threat to block the hedged flows leaves $Q$ unchanged and need not be carried out (\cref{prop:counter}(b)). The design also moves the burden. Washington would act on the flows itself if The Hague did not concede, and on our reading an order to stop servicing carries a Chinese counter-threat after Decrees 834 and 835 \citep{statecouncil2026supplychain, reedsmith2026decree}, so the concession has become more expensive for Europe (\cref{prop:counter}(a)(i)). The Dutch answer was voice: the prime minister and the trade minister took objections to Washington in April and June 2026, and the cabinet told parliament it opposed the bill's extraterritorial reach \citep{nltimes2026pushback, vangerven2026dutch, tk2026match}. The Hague did not threaten to stop aligning. With $Q<0$, voice and the hold-up of \cref{sec:holdup} are all a holder has.

China's share of ASML's system sales rose from 14 percent in the second quarter of 2026 to 23 percent in the third \citep{asml2026q2call, asml2026q3}. The model has no inventories, so a stockpiling reading is not one of its implications. A rise in Dutch licensing of legacy tools, which would be a larger tolerated hedge, is. Licensing decisions and order dates would separate the two.

\subsection{Comparison cases}\label{sec:cases}

\Cref{tab:cases} sets the Dutch record beside six contemporary and three historical cases, coded with the rules of \cref{sec:coding}. The equipment holders in Japan, South Korea and Germany face the structure of the model. Taiwan and the United Arab Emirates fall outside the scope conditions of \cref{sec:scope}: China threatens Taiwan rather than bidding for it, and the UAE has no treaty guarantee. The historical cases add variation in linkage and compensation.

\begin{table}[p]
\centering
\scriptsize
\setlength{\tabcolsep}{2.5pt}
\caption{US demands on holders, coded through the model}\label{tab:cases}
\begin{tabularx}{\textwidth}{@{}>{\raggedright\arraybackslash}p{1.75cm}*{6}{>{\raggedright\arraybackslash}X}@{}}
\toprule
Holder, period & US demand & Instrument threatened & Compensation & Rival bid or counter-threat & Outcome & Reading in the model\\
\midrule
Netherlands (ASML), 2019--2026 & End EUV, then advanced and older immersion DUV and metrology sales to China; from 2026, servicing & Flow jurisdiction: zero de minimis (2023), SME FDP rule (2024), MATCH Act (proposed) & None; Supplement No.~4 status, a threat left unexercised & China's market, 26--36\% of sales (2023--2025); Nexperia ban executed (2025) & Aligned on sales with lags; kept the licensing pen; servicing in China continued into 2026 & $Q\le0$ through jurisdiction (\cref{rem:firm}); hedge tolerated (\cref{cor:rationing}); security channel untested\\
\addlinespace
Japan (Tokyo Electron), 2023--2024 & Control advanced equipment; in 2024, restrict sales and servicing in China & Flow jurisdiction: FDP rule not used but not ruled out & None found; Supplement No.~4 status & China threatened ``severe economic retaliation''; Toyota warned of lost access to critical minerals & 23 equipment types licensed (2023), ``not aligning''; no new measures (2024); servicing continued into 2026 & Counter-threat raised the cost of the servicing concession; threat not executed (\cref{prop:counter}(a))\\
\addlinespace
South Korea (Samsung, SK hynix fabs in China), 2025 & No expansion or upgrade of Chinese fabs & Flow jurisdiction: VEU status revoked & None; annual licences to operate (2026); trade deal silent on controls & None documented & Accepted; fabs operate under licence & Hedge policed by flow jurisdiction; threat rebuilt as licences, not executed (\cref{prop:counter}(b))\\
\addlinespace
Germany (Zeiss, Trumpf), 2023--2024 & Join equipment controls; ``primary target is Zeiss'' & None documented & None & None documented & Reportedly declined; slower licensing of laser exports to China & Demand of low marginal value while Dutch licences govern the tools that carry Zeiss optics (\cref{cor:rationing})\\
\addlinespace
Taiwan, 2025 & Half of US-used chips made in the US (``50-50'') & Floated publicly while tariff talks were under way; no explicit threat found & None & China a threat, not a bidder & Rejected publicly & Relocation demand with high enforceability (\cref{cor:control}(c)); outside scope\\
\addlinespace
UAE, 2025--2026 & Align security regulations, prevent diversion, invest in US data centres & Conditional linkage: chip access & Client pays in investment; receives chip approvals and Country Group A:5 & Not documented in our sources & Licence-free access (July 2026), citing defence partnership and Operation Epic Fury & Roles reversed: Washington holds the chokepoint; explicit security-for-access reward; outside scope\\
\addlinespace
West Germany, 1961--1967 (offsets) & Offset the foreign-exchange cost of US troops & Conditional linkage: troop levels, explicit in 1966 & Holder paid: arms purchases, \$500m in US securities, gold pledge; US placed orders in the UK & German option to convert dollars into gold, given up & Redeployment of up to 35,000 US personnel agreed (1967); Erhard fell (1966) & Linkage extracting sunk payments; repeated renegotiation (\cref{prop:slack}(d)); domestic cost outside the model\\
\addlinespace
United Kingdom, 1957 (CoCom China differential) & Keep stricter controls on China than on the Soviet bloc & None; Washington chose no sanctions & None; US-backed compromise failed & China trade; 12 of 16 members for abolition & Differential abolished; US kept its own embargo & Secondary demand left without linkage (\cref{cor:rationing})\\
\addlinespace
West Germany, 1975--1977 (Brazil nuclear deal) & Defer or stop enrichment and reprocessing transfers & High-level pressure (``major crisis''); no explicit linkage, though Bonn feared for the US commitment & Offered to Brazil (fuel guarantee), not to Bonn & Brazil as buyer; Bonn cited its contract and the NPT & Contract kept; informal restraint on future transfers by 1978 & Concession on the forward margin; sunk deal kept\\
\addlinespace
Germany and Spain, 2026 (posture, access) & Germany: none stated; cut followed Merz's criticism of US strategy on Iran. Spain: base access, defence spending & Posture cut executed (Germany); trade threats (Spain) & None & None & About 5,000 troops ordered out of Germany; Spain kept its refusal & Executed instrument is spent (\cref{sec:conditionality}); access denial as hold-up (\cref{sec:holdup}); no link to chips\\
\bottomrule
\end{tabularx}
\par\smallskip
\begin{minipage}{\textwidth}\scriptsize
\emph{Sources by row.} Netherlands: \cref{sec:background} and \citet{bis2024fdpr, asml2026annual, zakaria2026match, baumgartner2026match, taylorwessing2025nexperia}. Japan: \citet{allen2023translation, hoganlovells2023japan, asiafinancial2024japan, taipeitimes2024retaliation, bis2024fdpr, zakaria2026match}. South Korea: \citet{bis2025veu, afs2025veu, reuters2025koreanlicence, census2026fablicense, whitehouse2025korea}. Germany: \citet{barkin2024zeiss, burger2023customs, kahle2023trumpf}. Taiwan: \citet{younger2025fiftyfifty, focustaiwan2025fiftyfifty}. UAE: \citet{whitehouse2025uae, commerce2025gulfchips, bis2026uae}. Offsets: \citet{frus1995foreignecon, frus1995westerneurope, frus1999germany, bundesbank1967blessing}. United Kingdom: \citet{frus1989economicdefense, hansard1957china}. Brazil deal: \citet{frus2015southamerica, frus2014westerneurope, frus2015nonproliferation}. Germany and Spain: \citet{martinez2026germany, roque2026spain, omullony2026ankara, rte2025nato}. ``None documented'' means that the sources we consulted record none.
\end{minipage}
\end{table}

\paragraph{What the comparison shows.} No holder in the table was paid in cash for aligning on a chokepoint. The currency is status (Supplement No.~4, validated end-user status, Country Group A:5), which BIS grants and withdraws by rule. Under our coding such statuses are exit costs, the pattern \cref{cor:master} leads us to expect from holders with negative premia. Where money moved, it moved from the client to the patron (the offsets, the UAE) or to a third party: in 1967 Washington placed orders in Britain to close the trilateral deal, and in 1977 it offered Brazil, not Bonn, a fuel guarantee \citep{frus1995westerneurope, frus2015nonproliferation}.

Resistance clusters on the installed base and on deals already signed. Tokyo and The Hague kept servicing Chinese fabs, Bonn kept its contract with Brazil while offering restraint on future transfers, and the Korean fabs kept operating under licence. In the model these concessions cost the holder much relative to the bite that remains, and in Japan a Chinese counter-threat sat on exactly the concession Washington sought. A cost-minimizing coercer that found enforcement against servicing expensive would also tolerate it, so this pattern fits our account without separating it from that rival.

Explicit links between an ally's conduct on one issue and US treatment on another appear in three places: Johnson's link between troop levels and offset payments in 1966, the posture and trade measures against Germany and Spain in 2026, and the BIS rule of July 2026 that cited the UAE's role in Operation Epic Fury. None of them concerns a European holder's chip controls. Washington uses the security channel, but the record does not show it using that channel on ASML.

\paragraph{Where the comparison strains the model.} The offsets fit only partly. Washington linked troop levels to German payments, and the payments were sunk concessions, which escape the slack cap (\cref{prop:slack}(b)). Renegotiation recurred from the 1961 memorandum through the informal settlement of 1964 to the trilateral talks of 1966--67 \citep{frus1995foreignecon, frus1999germany, frus1995westerneurope}, as \cref{prop:slack}(d) leads us to expect of a patron that cannot commit to stop demanding, though FRUS does not show Bonn reasoning that way. The redeployment of up to 35,000 personnel agreed in 1967 strains the model, since a complete-information patron does not both collect and cut. We see two readings. If the troops were at Washington's ex post optimum, the linkage was a costly conditional threat and the partial cut is unexplained. If they were above it, as balance-of-payments pressure and the Mansfield resolution suggest, keeping them was a transfer to Bonn, the offsets were its price, and Washington cut where cutting saved money \citep{frus1995westerneurope, gavin2004gold, zimmermann2002money}. FRUS does not let us choose. Erhard's fall is outside the model. McCloy recorded Brandt calling the package approach ``a major mistake'' that hastened it \citep{frus1999germany}. Our holder bears a sovereignty cost only when it transfers control, and the case suggests that other concessions to the patron carry a domestic cost too.

The United Kingdom in 1957 shows a secondary demand that Washington chose not to back with linkage. Eisenhower conceded privately that there was ``very little of profit in the matter'', Macmillan answered that ``We live by exports'', and in September the National Security Council decided that no sanctions would be applied against allies for eliminating the special China controls \citep{frus1989economicdefense, mastanduno1992economic}. A demand worth little to Washington and costly to the holder is the first to go under \cref{cor:rationing}. The strain is that the holder did not act alone: twelve of sixteen members backed abolition \citep{frus1989economicdefense}, a collective exit that our single-holder model does not represent. The Brazil deal adds a different strain. German officials feared for the US security commitment, although in the FRUS record Washington did not tie its pressure to it \citep{frus2014westerneurope, frus2015nonproliferation}. Conditionality that exists only in the holder's beliefs has no place in our complete-information model. And \citet{gray2012commercial} reads Bonn's resistance as commercial rivalry, against the finding of \citet{kroenig2009exporting} that sensitive nuclear transfers follow strategic motives.

\subsection{Designs that could test the model}\label{sec:tests}

\paragraph{An event dataset.} The unit would be a US demand on a holder of a chipmaking-equipment chokepoint, 2019--2026: the Netherlands, Japan, South Korea and Germany. Each demand would be coded, from the Federal Register, national gazettes, parliamentary records and firm filings, for the instrument threatened (\cref{sec:coding}), the holder's response, any compensation and any Chinese response. China's share of the firm's sales and credible Chinese counter-threats proxy for the premium. The model implies that conditional linkage against secondary demands appears when a holder has slack and fades as its premium rises toward zero, while the use of flow jurisdiction does not track the premium (\cref{prop:counter}(b)). It also implies that concessions in one issue area crowd out demands in another within the same period (\cref{cor:rationing}). Cost minimization implies neither.

\paragraph{An elite survey experiment.} Conditionality is a property of US behaviour, so the respondents should be US foreign-policy elites. A vignette would describe an ally holding a chokepoint and vary whether it complied with US export controls on China, refused, or complied on sales while servicing Chinese fabs. The gap in support for defending the ally, and for keeping US forces and intelligence sharing there, between compliant and noncompliant allies estimates the conditionality of the core guarantee and of posture. Further arms would vary whether the president had publicly tied protection to compliance (the stake of \cref{prop:conditionality}) and whether the ally hosts bases Washington uses (the sign in part (d)).

\paragraph{Process tracing.} Requests under the Dutch Open Government Act (Woo) for ministerial records of export-control talks with US counterparts from 2019 to 2026 would show whether security commitments, posture or the F-35 programme were ever raised, which is the evidence the security channel needs. The same records would date Dutch drafting against US threats, which separates coercion from shared interest better than outcomes do, and show what The Hague expected from China before 30 September 2025.

\paragraph{The Busan expiry.} The US suspension of the Affiliates Rule ran for one year from 10 November 2025, and China's suspension of its October 2025 controls ran to 10 November 2026 \citep{whitehouse2025busan, mofcom2025pause}. In September 2026 the two governments agreed to extend the arrangement that paused them to 10 January 2027, but no legal act extending the suspensions themselves had appeared by late October \citep{mofcom2026extension}. If both lapse, the rule behind the Nexperia demand and China's extraterritorial rare-earth controls return together \citep{mayerbrown2025rareearth}, and \cref{prop:counter}(a) leads us to expect Washington to move its demands on The Hague away from concessions that carry credible Chinese counter-threats, toward demands China cannot reach or flow jurisdiction it exercises itself. A second executed Chinese counter-threat against a Dutch step would be hard to square with misjudgment, which learning should correct, and the crowding-out reading would then require us to find the demand it displaced.
